% sec-00-cover-letter.tex - IND Cover Letter (final)
\section*{Cover Letter}
\addcontentsline{toc}{section}{Cover Letter}

\noindent ChemicalQDevice \\
San Diego, California \\
July 1, 2026 \\[1em]
Food and Drug Administration \\
Center for Drug Evaluation and Research \\
Division of Oncology 1 \\
5901 to 5913 Ammendale Road \\
Beltsville, Maryland 20705 to 1266 \\[1em]

\noindent\textbf{Re: Initial Investigational New Drug Application, Serial 0000 - Daraxonrasib (RMC-6236), Phase 1, LLM-Directed Robotic Whipple in KRAS-Mutated PDAC} \\[1em]

\noindent Dear Division of Oncology 1 Reviewer:

On behalf of ChemicalQDevice, and in my capacity as sponsor-investigator, I am submitting this initial Investigational New Drug Application (IND v1.0) under 21 CFR \S 312.23 for a Phase 1, first-in-human, prospective open-label single-arm study of daraxonrasib (RMC-6236), an oral once-daily RAS(ON) multi-selective (pan-RAS) small-molecule inhibitor, administered in the perioperative setting of a large language model (LLM) directed eight-arm robotic pancreaticoduodenectomy (Whipple) for patients with resectable or borderline-resectable, KRAS G12 mutated pancreatic ductal adenocarcinoma (PDAC) \cite{daraxwhipple,onpremwhippl}. This application is assigned serial number 0000 as the initial submission and is transmitted with FDA Form 1571 as the cover sheet for the submission and the sponsor agreement, together with FDA Form 3674 (clinicaltrials.gov certification) and FDA Form 1572 \cite{cfr312paiuni}. The mechanistic premise is well established: daraxonrasib binds GTP-bound active mutant RAS in complex with cyclophilin A and thereby suppresses downstream effector signaling, and the agent received FDA Breakthrough Therapy designation in June 2025 for previously treated metastatic PDAC bearing KRAS G12 alterations \cite{daraxwhipple,pdac060s2030}.

This submission is combined in nature. It is concurrently an IND under 21 CFR Part 312 for the investigational drug daraxonrasib and an Investigational Device Exemption (IDE) under 21 CFR Part 812 for the significant-risk, Class II collaborative on-premises robotic surgical device that directs the operative procedure as an early-feasibility study \cite{onpremwhippl,paiautospons}. The LLM functions strictly as a hash-pinned, second-opinion advisory oracle isolated outside the vendor kinematic stack; full device autonomy is prohibited in accordance with 21 CFR \S 312.21(e), and a verification-before-generation ten-gate suite governs every advisory output (ACCEPT, BLOCK, or ESCALATE) consistent with H.R. 9510, the Verification Before Generation in Physical AI Oncology Trials Act of 2026 \cite{congress9510,clintrustpai}. Dual regulatory pathways are presented in an integrated dossier so that the drug and device safety considerations are reviewed together rather than in isolation, and so that any drug-device interaction is visible to a single reviewing team.

We respectfully request the standard 30-day safety review period under 21 CFR \S 312.40, after which, absent a clinical hold under 21 CFR \S 312.42, the sponsor-investigator intends to initiate the study at a single academic site \cite{phase1guidance}. Dose-escalation design follows a 3+3 schema across three oral once-daily dose levels (DL1 160 mg, DL2 220 mg, DL3 300 mg, the RASolute 302 recommended Phase 2 dose) with a 28-day dose-limiting toxicity window, enrolling up to 18 treated subjects from approximately 36 screened, with staggered sentinel enrollment and accrual target of roughly 1-2 subjects per month \cite{daraxwhipple}.

\needspace{6\baselineskip}
\noindent\textbf{Study at a glance.} Table C.1 summarizes the principal design parameters so that the reviewer has the operative facts of the investigation on the first page of the submission.

\needspace{10\baselineskip}
\noindent\begin{tabularx}{\textwidth}{L{3.7cm} Y}
\toprule
\textbf{Parameter} & \textbf{Specification} \\
\midrule
Investigational drug & Daraxonrasib (RMC-6236), oral RAS(ON) multi-selective (pan-RAS) small-molecule inhibitor, once daily \\
Investigational device & On-premises LLM-directed eight-arm collaborative robotic Whipple (8 arms $\times$ 7 degrees of freedom = 56 degrees of freedom; 640 sensor channels), Class II, significant-risk, early-feasibility \\
Indication & Resectable or borderline-resectable PDAC with documented KRAS G12 (G12D, G12V, or G12R); ECOG 0 to 1; Whipple candidate \\
Phase and design & Phase 1, first-in-human, prospective open-label single-arm; combined IND (21 CFR Part 312) and IDE (21 CFR Part 812) \\
Dose levels & DL1 160 mg, DL2 220 mg, DL3 300 mg once daily; 3+3 escalation; 28-day dose-limiting toxicity window \\
Sample size & Up to 18 treated subjects (approximately 36 screened); 1 academic site; 1 to 2 enrollments per month \\
Co-primary endpoints & Device and procedure SAE incidence through 30 days; MTD and RP2D of daraxonrasib; technical feasibility (task completion without unplanned conversion attributable to device malfunction or unsafe behavior) \\
LLM role & Hash-pinned second-opinion advisory oracle only; full autonomy prohibited (21 CFR \S 312.21(e)); 10 gate verify before generation suite \\
\bottomrule
\end{tabularx}
\vspace{-0.55cm}
\figcaption{Table C.1. Parameters, Phase 1 LLM-directed robotic Whipple: daraxonrasib in KRAS G12 mutated PDAC.}

\clearpage

\needspace{6\baselineskip}
\noindent\textbf{Contents enclosed.} The following components are submitted in the ReGARDD IND order, with the Cover Letter and FDA Forms 1571 and 3674 preceding the generated Table of Contents:

\needspace{12\baselineskip}
\noindent\begin{tabularx}{\textwidth}{L{1.5cm} Y}
\toprule
\textbf{Item} & \textbf{Component} \\
\midrule
1 & FDA Forms 1571 and 3674 (3674 Box B selected; PHS Act 402(j) registration requirements addressed) \\
2 & Table of Contents (generated) \\
3 & Introduction, including drug name, pharmacological class, structure, formulation, route, and duration \\
4 & General Investigational Plan, including the 3+3 escalation rationale and first-year plan \\
5 & Investigator Brochure \\
6 & Proposed Clinical Research (study protocol v1.0.0, informed consent, Form 1572, financial disclosures) \\
7 & Chemistry, Manufacturing, and Control information (drug substance, drug product, placebo, labeling, environmental assessment) \\
8 & Pharmacology and Toxicology information, including the integrated computational evidence base \\
9 & Human Experience (RASolute 302, RASolve 301, Breakthrough Therapy designation) \\
10 & Additional Information (drug dependence, radioactive drugs, pediatric studies) \\
11 & Relevant Information and References with Back Matter \\
\bottomrule
\end{tabularx}
\vspace{-0.5cm}
\figcaption{Table C.2. Enclosed components of the initial IND, in ReGARDD order, with the Cover Letter and FDA Forms.}

\noindent\textbf{Regulatory basis for each enclosure.} Each component answers a specific content requirement of 21 CFR \S 312.23(a), and Table C.3 maps the enclosure to its codified citation so that the reviewer can confirm completeness against the regulation directly. The mapping also makes explicit where the combined drug and device structure adds content beyond the conventional drug-only IND, in particular the integrated Physical AI safety, governance, and verification material that supports the IDE arm \cite{cfr312paiuni,paiautospons}.

\needspace{14\baselineskip}
\noindent\begin{tabularx}{\textwidth}{L{6.3cm} L{3.1cm} Y}
\toprule
\textbf{IND content element} & \textbf{21 CFR citation} & \textbf{Where addressed} \\
\midrule
Cover sheet and sponsor agreement & \S 312.23(a)(1) & FDA Form 1571, serial 0000 \\
Table of contents & \S 312.23(a)(2) & Generated Table of Contents \\
Introductory statement and general investigational plan & \S 312.23(a)(3) & Introduction and General Investigational Plan \\
Investigator brochure & \S 312.23(a)(5) & Investigator Brochure \\
Protocols & \S 312.23(a)(6) & Proposed Clinical Research, v1.0.0 \\
Chemistry, manufacturing, and control & \S 312.23(a)(7) & CMC information \\
Pharmacology and toxicology & \S 312.23(a)(8) & Pharmacology and Toxicology, with computational evidence base \\
Previous human experience & \S 312.23(a)(9) & Previous Human Experience \\
Additional information & \S 312.23(a)(10) & Additional Information \\
Significant-risk device and safety reporting & 21 CFR Part 812; \S 312.32; \S 312.404 & IDE arm; integrated Physical AI safety reporting \\
\bottomrule
\end{tabularx}
\vspace{-0.3cm}
\figcaption{Table C.3. Cross-reference of each enclosed IND component to its governing provision of 21 CFR\\ \S 312.23(a) and to the supporting device and safety regulations engaged by the combined submission.}

\clearpage

\noindent\textbf{Sponsor-investigator and contact.} The sponsor-investigator is Kevin Kawchak, Chief Executive Officer, ChemicalQDevice, San Diego, California, who may be reached at kevink@chemicalqdevice.com. The sponsor-investigator's financial interest is disclosed under 21 CFR Part 54 and mitigated by a three-tier independent oversight structure comprising a Data Safety Monitoring Board with halt authority at cohort boundaries and on triggered review, an Independent Safety Monitor with real-time per-procedure stop authority, and a Physical AI Safety Review Committee operating on a cadence of no more than 90 days \cite{paisitedocpk}. Binding halt rules require Data Safety Monitoring Board review and possible halt upon any device or procedure serious adverse event in at least one of three subjects within a cohort, or upon two device-related deaths at any time, providing an explicit and conservative stopping boundary for this first-in-human Physical AI investigation.

\needspace{6\baselineskip}
\noindent\textbf{Safety reporting commitments.} The sponsor-investigator will comply with the expedited reporting timelines of 21 CFR \S 312.32 and with the Physical AI reporting framework engaged by the device arm. Table C.4 states controlling clocks that will govern the conduct of the study from first dose onward \cite{cfr312paiuni}.

\needspace{12\baselineskip}
\noindent\begin{tabularx}{\textwidth}{L{7cm} L{3.6cm} Y}
\toprule
\textbf{Reportable event} & \textbf{Clock} & \textbf{Basis} \\
\midrule
Fatal or life-threatening suspected adverse reaction & 7 calendar days & 21 CFR \S 312.32(c)(1)(i)(A) \\
Other serious and unexpected suspected adverse reaction & 15 calendar days & 21 CFR \S 312.32(c)(1)(i)(B) \\
Serious Physical AI adverse event & 7 or 15 calendar days & 21 CFR \S 312.32(f), (g) \\
Cybersecurity incident with potential for patient harm & 7 calendar days & 21 CFR \S 312.32(g) \\
Cybersecurity incident without patient-harm potential & 15 calendar days & 21 CFR \S 312.32(g) \\
Sustained AI model degradation, sim-to-real divergence beyond twice tolerance, or digital-twin discrepancy & 15 calendar days & 21 CFR \S 312.32(g) \\
\bottomrule
\end{tabularx}
\vspace{-0.35cm}
\figcaption{Table C.4. Expedited safety reporting clocks the sponsor-investigator commits to under 21 CFR \S \\312.32, including the six Physical AI reporting triggers engaged by the investigational device.}

\noindent\textbf{AI-assisted assembly and real-time monitorability.} This IND was assembled under a single master prompt that executed a four-stage mermaid to draft to full to final build, with each generated file committed and pushed to one continuously updated GitHub pull request in real time, so that the complete authoring history of the submission is human-monitorable as it is produced and is reproducible from a hash-pinned deterministic commit (repository v4.3.0) \cite{oncotrialllm,natlpaiplatf}. The same verification-before-generation discipline applied to clinical advisory outputs governs document generation: each figure is reproduced exactly from a grayscale Mermaid source, each numeric value is grounded in the protocol and the computational evidence base, and each codified citation is checked against the regulation \cite{clintrustpai,llmcomp16fda}. Figure 1 shows the single-prompt build pipeline and Figure 2 shows the ReGARDD Table-of-Contents architecture that maps each numbered IND section to exactly one source file. The repository is available at \url{https://github.com/kevinkawchak/physical-ai-oncology-trials/tree/main/trial-ind}.

\needspace{8\baselineskip}
\begin{mermaidfig}
\node[mmgoal] (mp) at (0,0) {Master prompt\\prompts/prompt-ind.md};
\node[mmproc] (pa) at (4.6,0) {Process A\\write sub-prompts 1 to 4};
\node[mmin] (s1) at (0,-2.4) {Stage 1 mermaid\\22 grayscale figures};
\node[mmin] (s2) at (4.6,-2.4) {Stage 2 draft-ind\\12 sec tex scaffold\\+ instructions};
\node[mmaccent] (s3) at (9.2,-2.4) {Stage 3 full-ind\\full prose + 22 TikZ\\+ full-width tables};
\node[mmgoal] (s4) at (9.2,-5.0) {Stage 4 final-ind\\senior-author polish\\+ Overleaf zip};
\node[mmdark] (rel) at (4.6,-5.0) {Release v4.3.0\\README + CHANGELOG\\+ output-ind};
\node[mmctx] (gh) at (0,-5.0) {GitHub branch\\real-time auto-commit\\and auto-PR};
\draw[mmedge] (mp) -- (pa);
\draw[mmedge] (pa) to[out=-90,in=90,looseness=0.6] (s1.north);
\draw[mmedge] (s1) -- (s2);
\draw[mmedge] (s2) -- (s3);
\draw[mmedge] (s3) -- (s4);
\draw[mmedge] (s4) -- (rel);
\draw[mmedged] (s1) -- (gh);
\draw[mmedged] (s2.south) to (gh.north);
\draw[mmedged] (s3.south west) to (gh.north);
\draw[mmedged] (s4.south) to[out=-110,in=-20,looseness=0.3] (gh.south);
\draw[mmedged] (gh.north) to[out=70,in=120,looseness=0.2] (s4.north);
\end{mermaidfig}
\vspace{-0.75cm}
\figcaption{Figure 1. The single-prompt mermaid to draft to full to final build pipeline for the Phase 1 PDAC IND. One master prompt drives Process A (write the four sub-prompts) and then Process B (execute them as Stages 1 to 4). Dashed edges show that every generated file is committed and pushed to one continuously updated pull request in real time, and the feedback edge shows author review returning into the final stage, all without manual intervention.}

\needspace{8\baselineskip}
\begin{mermaidfig}
\node[mmaccent] (cl) at (0,0) {Cover Letter\\sec-00};
\node[mmproc] (f) at (0,-1.8) {1. FDA Forms 1571 / 3674\\sec-01};
\node[mmctx] (toc) at (0,-3.6) {2. Table of Contents\\generated};
\node[mmin] (i) at (0,-5.4) {3. Introduction\\sec-02 (3.1 to 3.5)};
\node[mmin] (gip) at (4.6,0) {4. General Investigational Plan\\sec-03 (4.1 to 4.7)};
\node[mmin] (ib) at (4.6,-1.8) {5. Investigator Brochure\\sec-04};
\node[mmin] (pcr) at (4.6,-3.6) {6. Proposed Clinical Research\\sec-05 (6.1 to 6.3)};
\node[mmin] (cmc) at (4.6,-5.4) {7. Chemistry, Manufacturing, Control\\sec-06 (7.1 to 7.2)};
\node[mmin] (pt) at (9.2,0) {8. Pharmacology / Toxicology\\sec-07 (8.1)};
\node[mmin] (phe) at (9.2,-1.8) {9. Previous Human Experience\\sec-08 (9.1 to 9.4)};
\node[mmin] (ai) at (9.2,-3.6) {10. Additional Information\\sec-09 (10.1 to 10.5)};
\node[mmin] (ri) at (9.2,-5.4) {11. Relevant Information\\sec-10};
\node[mmgoal] (bm) at (9.2,1.6) {References + Back Matter\\sec-11};
\node[mmgoal] (main) at (4.6,1.6) {main.tex\\assembles all};
\draw[mmedge] (cl) -- (f);
\draw[mmedge] (f) -- (toc);
\draw[mmedge] (toc) -- (i);
\draw[mmedge] (i.east) to[out=0,in=180,looseness=0.6] (gip.west);
\draw[mmedge] (gip) -- (ib);
\draw[mmedge] (ib) -- (pcr);
\draw[mmedge] (pcr) -- (cmc);
\draw[mmedge] (cmc.east) to[out=0,in=180,looseness=0.6] (pt.west);
\draw[mmedge] (pt) -- (phe);
\draw[mmedge] (phe) -- (ai);
\draw[mmedge] (ai) -- (ri);
\draw[mmedge] (ri.east) to[out=0,in=0,looseness=0.5] (bm.east);
\draw[mmedged] (main.south west) to[out=160,in=0,looseness=0.7] (cl.east);
\draw[mmedged] (main.east) -- (bm.west);
\draw[mmedged] (main.south west) to[out=180,in=0,looseness=0.6] (toc.east);
\end{mermaidfig}
\vspace{-0.75cm}
\figcaption{Figure 2. The ReGARDD IND Table of Contents reproduced as the file architecture of the build. The\\Cover Letter and the FDA Forms 1571 and 3674 precede the generated Table of Contents; the eleven\\numbered IND sections and the References and Back Matter each become exactly one sections/sec tex\\file assembled by main.tex, so the document structure and the repository structure match one-to-one.}

\clearpage

\noindent\textbf{Scientific and computational support.} The pharmacology and toxicology section is supported by an integrated computational evidence base developed by the sponsor-investigator across 2025 and 2026, including a 60-second PDAC Whipple and daraxonrasib simulation, an AI digital-twin PDAC simulation, a quantitative systems pharmacology metastatic-PDAC model with verification, validation, and uncertainty quantification, and a 100,000-patient in silico Phase 3 five-arm PDAC study performed in triplicate \cite{qspmetpancre,fdadigtwinpc,chatgpt100kp,pdacdigtwinp}. These analyses do not substitute for the empirical safety review but provide quantitative priors that inform the dose levels, the perioperative pause-and-restart advisory, and the sim-to-real tolerances that the device must satisfy before any patient-facing robotic activity. The clinical rationale rests on the severity of the indication: PDAC carries a five-year overall survival below 13 percent and ranks third in United States and fourth in European cancer mortality, so even a first-in-human feasibility study is justified by the magnitude of unmet need and by the precision of the proposed estimation framework \cite{pdac060s2030,paipredict4x}.

\needspace{6\baselineskip}
\noindent\textbf{Statistical framing.} The study is descriptive and non-confirmatory in design, consistent with ICH E9 and TRIPOD+AI, with no formal power calculation; precision is reported through exact Clopper-Pearson 95 percent confidence intervals \cite{ichpaistduni,llmcomp16fda}. As a worked example grounded in the maximum sample size, an observed rate of 0 events in 18 treated subjects corresponds to a point estimate of 0 percent with an exact 95 percent interval of 0 to 18 percent, while a single event in 18 subjects corresponds to 5.6 percent with an interval of 0.1 to 27 percent; at the cohort level, 0 of 6 corresponds to an interval of 0 to 46 percent. These intervals make explicit the limits of inference at first-in-human scale and frame all supportive two-sided testing at alpha 0.05 without multiplicity adjustment \cite{ichpaistduni}. Informed consent will be obtained under 21 CFR Part 50 and ICH E6(R3), with a Physical AI opt-out under 21 CFR \S 312.60(f), explicit treatment of therapeutic misconception, and a plain-language explanation of staged autonomy, which for this Phase 1 is limited to clinician-controlled teleoperation and supervised operation only \cite{cfr050paiuni}.

\noindent We appreciate the Division's review and welcome any questions during the 30-day period. The sponsor-investigator is available for a teleconference at the Division's convenience and will respond promptly to any request for additional information.

\noindent Respectfully submitted, \\[1.0em]
Kevin Kawchak \\
Chief Executive Officer and Sponsor-Investigator \\
ChemicalQDevice, San Diego, California \\
kevink@chemicalqdevice.com \\
IND v1.0; repository v4.3.0
